% !TEX root = ./CV_Tramarin.tex

%%%%%%%%%%%%%%%%%%%%%%%%%
%%%%%%% RESEARCH PROGRAMS
%%%%%%%%%%%%%%%%%%%%%%%%%

\sectioniten{Direzione o partecipazione alle attività di un gruppo di ricerca caratterizzato da
	collaborazioni a livello nazionale o internazionale}{Direction or Participation of Research Groups, characterized by National or International Collaborations}


\cventryit{2021 - presente}{Co-Direzione del Gruppo di ricerca Internazionale ``Intel Wireless TSN''}
{}{\newline \emph{Tematiche di ricerca}: Sviluppo e valutazione delle prestazioni di sistemi di misura distribuiti innovativi e ad alte prestazioni tramite basati su reti wireless TSN.}
{\newline \emph{Ruolo personale}: Co-fondatore del gruppo di ricerca, Responsabile dell'Unità di Modena}{ Unità di ricerca:
	\begin{itemize}
		\item Intel Labs (Oregon, USA)
		\item Università di Modena e Reggio Emilia
		\item Università di Padova
		\item Consiglio Nazionale delle Ricerche
	\end{itemize}
	}

\cventryen{2021 - present}{Co-Director of the International Research Group ``Intel Wireless TSN''}
	{}{\newline \emph{Research topics}: Development and performance assessment of innovative and high-performance distributed measurement systems based on wireless TSN networks.}
	{\newline \emph{Personal role}: Co-Founder of the research group, Director of the Modena unit} { Research units:
	\begin{itemize}
		\item Intel Labs (Oregon, USA)
		\item Università di Modena e Reggio Emilia
		\item Università di Padova
		\item Consiglio Nazionale delle Ricerche
	\end{itemize}
	}

\cventryit{2020 - presente}{Co-Direzione del gruppo di ricerca IoT in Factory Automation}
	{}{\newline \emph{Tematiche di ricerca}: Ricerca sulle tecnologie di misura e networking nel settore dell'automazione industriale e di fabbrica, con un approccio interdisciplinare basato su misurazioni, controllo e informatica.}
	{\newline \emph{Ruolo personale}: Co-fondatore del gruppo di ricerca, Responsabile dell'Unità di Modena}{Partecipanti: Gruppo di ricerca IoT (Dipartimento di Ingegneria, Unimore), Università di Padova, Consiglio Nazionale delle Ricerche (CNR-IEIIT)}
	{}%Ruolo personale: Fondatore del gruppo di ricerca, Direttore dell'unità di Modena}%{}

\cventryen{2020 - present}{Co-Director of the IoT in Factory Automation research group}
		{}{\newline \emph{Research topics}: Research on measurement and networking technologies in the field of industrial and factory automation, with an interdisciplinary approach based on measurements, control, and computer science.}
		{\newline \emph{Personal role}: Co-Founder of the research group, Director of the Modena unit}{Participants: IoT research group (Dept. Engineering, Unimore), University of Padova, National Research Council (CNR-IEIIT)}
		{}%Personal role: Founder of the research group, Director of the Modena unit}%{}

\cventryit{2021 - 2023}{Direzione e responsabilità scientifica del gruppo di ricerca SHeMS}
	{\newline costituito per il progetto ``Smart Healthcare Modular System for Personalized Remote Health Monitoring'', Call no. FARMO --- 1/12/2021\newline \emph{Tematiche di ricerca}: Progettazione, sviluppo e test di un nuovo sistema sanitario intelligente modulare per il monitoraggio remoto personalizzato della salute}
	{\newline Partecipanti: Gruppo di ricerca IoT (Dipartimento di Ingegneria, Unimore) e Gruppo di ricerca in Cardiologia (Policlinico Unimore)}{}
	{}

\cventryit{}{Tutor scientifico di un assegno di ricerca}{}{Nel ruolo di Responsabile scientifico del grppo di ricerca SHeMS, il cadidato è stato anche Tutor scientifico di un assegno di ricerca (Bando 2021-adrj064), della durata di 12 mesi. L'assegno di ricerca è stato assegnato alla Dott.ssa Valentina Di Pinto}{}{}{}


\cventryen{2021 - 2023}{Director and Principal Investigator of the SHeMS research group}
	{\newline established for the project ``Smart Healthcare Modular System for Personalized Remote Health Monitoring'', Call no. FARMO --- 1/12/2021\newline \emph{Research topics}: Design, development, and testing of a novel smart healthcare modular system for personalized remote health monitoring}
	{\newline Participants: IoT research group (Dept. Engineering, Unimore) and Cardiology Research group (Polyclinic Unimore)}{}
	{}

% \cventryen{}{Scientific tutor of a research grant}{}{In the role of Scientific Director of the SHeMS research group, the candidate was also the Scientific tutor of a research grant (Call 2021-adrj064), lasting 12 months. The research grant was awarded to Dr. Valentina Di Pinto}{}{}{}

